\documentclass[12pt]{article}
\usepackage[utf8]{inputenc}
\usepackage[T1]{fontenc}
\usepackage{hyperref}
\title{Waata Language Resources for K-12 Curriculum Development}
\author{Compiled by Claude Sonnet 5.5 (claude-sonnet-5-5) for Andr\'as Kornai}
\date{}
\begin{document}
\maketitle

\section{Introduction}
Waata (ISO 639-3 ssn; also called Sanye) is described as a dialect of Southern Oromo spoken in Kenya. Search results give about 20,000 speakers (2019 census, per Wikipedia) and about 13,000 (Omniglot); the figures conflict. A Latin-script orthography exists, but its use is limited. Very little educational material was found, and nothing indicates that Waata is taught in schools.

\section{Primary Educational Resources}
No primer, textbook or school programme for Waata turned up in the searches. Curriculum work would have to start from the materials below and from new development with the community.

\section{Linguistic and Reference Materials}
\begin{itemize}
\item Bernd Heine, \emph{The Waata Dialect of Oromo: Grammatical Sketch and Vocabulary} (Language and Dialect Atlas of Kenya, vol.\ 4, Reimer, Berlin, 1981, 72 pp.). This is a short sketch with a vocabulary list, not a full dictionary. Bibliographic records: \url{https://glottolog.org/resource/reference/id/82381} and \url{https://quaibranly.fr/en/explore-collections/base/Bibliography/action/show/notice/000084582-the-waata-dialect-of-oromo-grammatical-sketch-and-vocabulary}
\item General description of the language: \url{https://en.wikipedia.org/wiki/Waata_language}
\item Writing system overview: \url{https://www.omniglot.com/writing/waata.htm}
\item Endangered Languages Project entry: \url{https://www.endangeredlanguages.com/lang/4787?hl=de}
\end{itemize}

\section{Religious and Community Resources}
According to Wikipedia, a Waata New Testament was released in July 2025 after a project overseen by Wycliffe Bible Translators and Bible Translation and Literacy, and Old Testament work has begun. The Jesus Film was launched in the language in 2018 (\url{https://en.wikipedia.org/wiki/Waata_language}). I did not find the scripture text itself in the search results, so its availability and orthography need checking directly with the translation organisations. No community organisation was identified.

\section{Digital and Audio Resources}
Global Recordings Network lists audio-visual Bible lessons and the Look, Listen \&amp; Live series for Waata: \url{https://globalrecordings.net/en/language/ssn}

\section{Curriculum Development Context}
The documentation is thin: one 1981 grammatical sketch, a recently translated New Testament, and gospel audio. A native-language curriculum would first need a community-agreed orthography, a lexicon and a corpus of everyday speech. The Oromo-related literature may help with grammar, but Waata must not be assumed identical to standard Oromo. No existing school use was found.

\end{document}